\section{A collision-free Jones query and full polynomial recovery}
\label{sec:faithful-jones}

The fixed-color queries above can certify a nontrivial Jones value, but equality
can hide a nonconstant polynomial.  The optional \code{potts-faithful} backend
removes this specialization ambiguity.  It chooses an input-dependent quadratic
integer and either decides whether the \emph{whole} normalized Jones polynomial
is one, or recovers that polynomial.  With resource caps disabled, both queries
have the general bit bound
\[
 \operatorname{poly}(n)\,2^{O(\sqrt n\log(n+1))}.
\]
This is subexponential in the crossing count.  It is a bound for Jones
computation, not for complete unknot recognition.  The implementation does not
assume that the Jones polynomial detects the unknot: a completed identity
result still proceeds to an independent recognition certificate.

This section combines elementary coefficient bounds and integer encoding with
the preceding maintained separator and Potts algorithms.  It makes no priority
claim for subexponential Jones computation.  In particular, efficient Jones
computation under graph-width restrictions predates this implementation; see
the literature review in \code{reports/17/research/literature\_landscape.md}.

\subsection{Uniform coefficient and degree bounds}

Write $V_D(t)$ for the normalized Jones polynomial of an $n$-crossing knot,
with $V_{\bigcirc}=1$, and put $W(x)=V_D(x^{-1})$ for $x=A^4$.
The bracket state sum has $2^n$ states.  Each state has at most $n+1$ circles:
changing one smoothing changes the circle count by one, and smoothing a
connected $n$-crossing projection cannot produce more than $n+1$ components.
The normalized circle factor is $(-A^2-A^{-2})^{c-1}$, whose coefficient
$\ell^1$ norm is $2^{c-1}$.  Writhe normalization only changes signs and
shifts degrees.  Thus
\[
 \|W\|_1=\|V_D\|_1\le 2^n 2^n=4^n=:C.
\]
The smoothing monomial has $A$-exponent of absolute value at most $n$;
circle factors contribute at most $2n$; writhe contributes at most $3n$.
Every resulting exponent is divisible by four for a knot, giving
$|\deg_x|\le 3n/2$.  We use the simpler integral bound $D=2n$.
The crossing-free knot is handled separately with polynomial one.

\subsection{Choosing a specialization that cannot hide nontriviality}

Choose
\[
 M=2^{4n+2}=4\cdot16^n,\qquad q=M+2,\qquad x^2-Mx+1=0,
\]
and consider the larger real root $x$.  Since $M\ge4$, this root satisfies
$x>M-1>C$.  Suppose that $W-1$ is a nonzero Laurent polynomial.  Multiplying
by a monomial produces an ordinary integer polynomial $P$ with nonzero
leading coefficient of absolute value at least one.  Its norm is at most
$C+1$, so the sum of absolute values of all its other coefficients is at most
$C$.  At a real argument $x>C$, the leading term has absolute value greater
than the sum of the lower terms.  Thus $P(x)\ne0$, and hence $W(x)\ne1$.

Consequently equality of the exact quadratic-ring pairs for the partition
function $Z$ and the unknot normalization $U$ is equivalent to
$V_D(t)=1$ as a polynomial.  The normalization is nonzero: its formula is a
signed monomial times a power of $x+1$.  No floating-point approximation or
probabilistic identity test is used.  The same leading-term argument applied
to a difference of two bounded polynomials proves injectivity, because
$M-1>2C$.

The root argument alone would permit a smaller base for identity testing.
We deliberately use a single larger base for both identity testing and the
coefficient recovery below, keeping their result semantics consistent.

\subsection{Recovering all coefficients from two integers}

The existing evaluator supplies $Z$ and $U$ as pairs representing elements
of $\mathbb{Z}[x]/(x^2-Mx+1)$.  First recover the normalized pair exactly.
If $U=c+dx$, its conjugate is $(c+Md)-dx$ and its norm is
\[
 N(U)=c^2+Mcd+d^2.
\]
Multiply $Z$ by this conjugate and divide both coordinates by $N(U)$.
The divisions have zero remainder because $Z/U=W(x)$ is a Laurent
polynomial with integer coefficients and $x^{-1}=M-x$ is integral.
The code checks the remainders, rather than rounding them.

To understand the coordinate encoding, introduce formal polynomials
$U_{-1}(M)=0$, $U_0(M)=1$, $U_1(M)=M$, and
$U_j(M)=M U_{j-1}(M)-U_{j-2}(M)$.  For $k\ge1$,
\[
 x^k=-U_{k-2}(M)+U_{k-1}(M)x,\qquad
 x^{-k}=U_k(M)-U_{k-1}(M)x.
\]
The first formula for $k=1$ uses $U_{-1}=0$.
Each $U_j$ has coefficient norm at most $2^j$, by induction on the
recurrence.  Therefore reduction of $W$ yields coordinate polynomials
$a(M),b(M)$ of degree at most $D$, each with every coefficient bounded by
\[
 C 2^D\le 4^n 2^{2n}=16^n=M/4.
\]
The computed integers $a(M),b(M)$ consequently determine these coordinate
polynomials uniquely by balanced base-$M$ expansion.  Repeatedly take the
residue modulo $M$, subtract $M$ when it is at least $M/2$, and divide the
remaining difference by $M$.  This handles negative coordinates as well as
positive ones.  Actual coefficients lie strictly inside the balanced digit
interval; there is no carry ambiguity.

Finally substitute the formal identity $M=x+x^{-1}$.  A decoded term
$cM^j x^e$, where $e\in\{0,1\}$, contributes
\[
 c\sum_{k=0}^j\binom jk x^{j-2k+e}.
\]
Aggregate equal exponents, remove zero coefficients, and negate the exponents
to return the standard variable $t=A^{-4}$.  The implementation independently
checks the final degree and coefficient bounds, and consistency with the
pair-equality identity result.  These checks detect inconsistent input to the
decoder; they do not replace the state-sum correctness argument.
There are $O(n)$ output coefficients of $O(n)$ bits each.  Balanced decoding,
exact normalization division and binomial expansion take polynomial bit time.

\subsection{Why exponentially many colors do not cause exponential work}

Applying the fixed-$q$ bound $q^f$ directly would lose the desired result.
The implementation instead stores \emph{canonical equality partitions} of
the $f$ active Tait vertices.  There are at most the Bell number $B_f$ such
partitions, independent of $q$.  Its extension loop considers each existing
color class and at most one new class, weighted by the number of unused
colors $q-k$.  It never enumerates $q$ colors.

The certified crossing order has frontier $w=O(\sqrt n)$, hence at most
$f\le w/2$ active faces of either chosen checkerboard shade.  As
$B_f\le\max(1,f^f)$, the number of represented states is
$2^{O(\sqrt n\log(n+1))}$.  Each crossing introduces at most two vertices;
the number of extensions from one state is at most $(n+3)^2$.
The original scalar coefficient bound is $O(n\log q)$ bits, now $O(n^2)$.
Key manipulation, multiplicities, arithmetic and the number of stages
therefore contribute only polynomial factors.

The adaptive preliminary scan needs separate accounting because $q$ is no
longer fixed.  Before its first threshold crossing it has at most $64n+1$
states.  One completed-stage overshoot costs at most
$(64n+1)(n+3)^2$ transitions.  If it takes the short-tail option with at most
$2\lceil\log_2(n+1)\rceil$ crossings remaining, a conservative work bound is
\[
 \operatorname{poly}(n)(n+3)^{O(\log(n+1))}
 =2^{O(\log^2(n+1))}.
\]
Thus the mode called \code{polynomial-tail} in the shared fixed-color policy
has only a \emph{quasi-polynomial} bound here.  That cost is still dominated
by $2^{O(\sqrt n\log(n+1))}$.  All other uncapped runs obtain the separator
certificate and continue or restart once.  Discarded work and a possible
state-cap interruption obey the same bounds.  The polynomial construction
and verification costs do not change the final bound.

\subsection{Interface, evidence and resource semantics}

\code{faithful\_potts\_exact} in \code{fastunknot/faithful\_jones.py}
returns the existing exact query fields together with
\code{polynomial\_identity}: the algorithm version, crossing count, whether
the polynomial is one, and the exponents defining the coefficient and
encoding bounds.  Its optional \code{include\_polynomial=True} returns
sorted Laurent coefficients in \code{jones\_polynomial}, with variable
\code{t} and signed hexadecimal coefficient strings.  Raw quadratic pairs
and the raw color count remain Python integers.

The command \code{python -B -m fastunknot jones FILE --backend potts-faithful}
requests full recovery.  Recognition through
\code{--jones-backend potts-faithful} requests identity testing only.
The color count is derived from the input; \code{--potts-colors} does not
override it.  The Python API and recognition defaults remain 4096 states and
200000 transitions; the standalone \code{jones} command has no local caps
unless supplied.  An exhausted allowance remains inconclusive.
The uncapped complexity statement
requires \code{max\_states=None} and \code{max\_transitions=None} and no
external deadline.  A restarted evaluation shares the original transition
allowance.  Global checks run during decoding too.  No incomplete identity
or polynomial is published if the query fails before completion.

Exact witnesses encode $q$ in \code{q\_hex} above 256 bits and retain the
existing integer \code{q} field for smaller values.  Large ring descriptions
refer symbolically to $q$.  This avoids Python's decimal-conversion limit
without changing process-global settings.  The factorized exact backend
uses the same ring description and witness conversion.

Identity equality has two deliberately separate consequences: it settles
the algebraic question $V_D=1$, but yields \code{INCONCLUSIVE} for the Jones
recognition filter.  The complete recognizer retains this stronger evidence
and continues to its existing exact fallback.  Neither this backend nor the
adaptive scheduler establishes a general subexponential recognition bound.

\subsection{Validation and measurement}

The complete maintained suite passes all 602 tests in 111.769 seconds;
\code{data/faithful-jones-integrated-tests.txt} preserves the run.
Seven new maintained tests compare 156 full polynomials against the independent
whole-cube Laurent oracle: both checkerboard shades and both mirrors of 35
random diagrams, three named fixtures, and the crossing-free knot.  Another
187 synthetic bounded Laurent polynomials check recovery at extreme signed
degrees and coefficient sizes, including arbitrary exact normalization.
Tests exercise a real certified restart, the precise shared transition limit,
the known $q=5$ weaving collision, recognition fallback after polynomial
identity, cancellation during decoding, zero budgets, CLI output, and
20001-bit color counts through both exact witness backends.

\code{benchmark\_faithful\_jones.py} measures fixed $q=6$ twice as an A/A
control, faithful identity, and full recovery in seven shuffled rounds after
an excluded warmup.  Every timing includes construction of a fresh diagram
and all order preparation.  Small inputs have independent cube checks;
completed larger recovered polynomials are checked by specialization back
to the fixed-$q$ result.  These arms deliver different guarantees, so their
timings describe the price of stronger information rather than speedups for
an identical task.  Capped results are retained as such.  The backend remains
optional, and these invariant timings do not claim faster full recognition.

\begin{center}\small
\begin{tabular}{@{}lrrrr@{}}
\toprule
Input & Fixed $q=6$ & A/A & Identity & Full polynomial \\
\midrule
Conway & 0.563 & 0.568 & 0.608 & 0.718 \\
Hard unknot 8 & 0.588 & 0.588 & 0.647 & 0.673 \\
Weaving 10 & 0.553 & 0.557 & 0.598 & 0.648 \\
Tree medial 254 & 17.981 & 18.105 & 125.268 & 239.190 \\
Grid 64, shuffled & 21.704 & 21.792 & 25.720 & 25.703 \\
Grid 100 & 21.873 & 21.874 & 38.744 & 38.943 \\
\bottomrule
\end{tabular}
\end{center}
Median milliseconds including setup; distinct output guarantees. A limit is a censored query.


All 13 measured inputs completed in every arm and round.  The Conway knot
costs 0.563\,ms at fixed $q=6$, 0.608\,ms for collision-free identity, and
0.718\,ms for full recovery.  The larger integer arithmetic is more costly
on the 254-crossing tree medial: 17.981\,ms, 125.268\,ms, and 239.190\,ms,
respectively.  Grid 100 rises from 21.873\,ms to 38.744\,ms for identity
and 38.943\,ms for recovery.  This cost is a reason to retain the inexpensive
fixed specializations and the current default filter; stronger asymptotic
information is not a measured per-input runtime improvement.
